\answer{ex:03:1}{$35\to17.5\mV$ при $g_E=g_L$; $56\to40\mV$ при $g_E=4g_L$, віднімання дало б $38.5\mV$: шунт ділить $g_E$ на $3$. Вікно звужується вдвічі ($\tau_{\text{еф}}$ падає у $2$ рази).}
\answer{ex:03:2}{$\operatorname{tr}=\frac{w_{EE}-1}{\tau_E}-\frac1{\tau_I}$, $\det=\frac{1-w_{EE}+w_{EI}w_{IE}}{\tau_E\tau_I}$; на межі $\omega=\sqrt{\det}$, $\tau_I=20\ms$, період $73\ms$.}
\answer{ex:03:3}{$d_c=\tau_E\arccos(-1/G)/\sqrt{G^2-1}$, період $2\pi\tau_E/\sqrt{G^2-1}\in(2d_c,4d_c)$. $G=2$: $d_c=12.1\ms$, період $36\ms$; $G=5$: $d_c=3.6\ms$, період $12.8\ms$ --- якраз смуга швидких ритмів.}
\answer{ex:03:4}{(а)~При $I_2=\beta I_1=2$ угору й $I_2=I_1/\beta=0.5$ униз, петля завширшки $1.5$. (б)~На $\tau\ln10/(\beta-1)=2.3\tau$ на кожен десятковий порядок $\varepsilon$.}
\answer{ex:03:5}{Три посередники, $d_k=5$, $10$, $15\ms$: заборони мають покрити $15\ms$, по $5\ms$ на кожну.}
\answer{ex:03:6}{$n+M+1=3+4+1=8$ нейронів. Два: \nt{GABA} $=[x+y+z\ge2]$, вихід $=[x+y+z-2\,\nt{GABA}\ge1]$.}
\answer{ex:03:7}{$\binom84=70<100\le\binom94=126$: $9$ посередників (теорема Шпернера). Без прямих зв'язків --- $100$: ранг $\beta(J-I)$ дорівнює $n$.}
\answer{ex:03:8}{$h_c=\frac1{2w_{EE}}+\frac{w_{EI}w_{IE}-w_{EE}}{4w_{EE}^2}=\frac7{18}\approx0.39$; моделювання дає зміну знака між $h=0.38$ і $0.40$.}
